\documentclass[journal]{IEEEtran}

\usepackage{amsmath,amssymb,amsthm}
\usepackage{array}
\usepackage{booktabs}
\usepackage{cite}
\usepackage{graphicx}
\usepackage{multirow}
\usepackage[ruled,vlined,linesnumbered]{algorithm2e}

\newcommand{\E}{\mathbb{E}}
\newcommand{\Prb}{\mathbb{P}}
\newcommand{\ind}{\mathbf{1}}
\newcommand{\cK}{\mathcal K}
\newcommand{\cC}{\mathcal C}
\newcommand{\cF}{\mathcal F}
\newcommand{\cP}{\mathcal P}
\newcommand{\cS}{\mathcal S}
\newcommand{\cW}{\mathcal W}
\newcommand{\FRT}{\mathrm{FRT}}
\newcommand{\MC}{\mathrm{MC}}
\newcommand{\sat}{\operatorname{sat}}

\newtheorem{theorem}{Theorem}
\newtheorem{lemma}{Lemma}
\newtheorem{proposition}{Proposition}

\title{Transition Aware Reliability Assessment of Hybrid AC/DC Distribution
Systems with Inverter Based Resources}

\author{Runze~Han,
        Tianyang~Zhao,~\IEEEmembership{Member,~IEEE,}
        and~Peng~Wang,~\IEEEmembership{Fellow,~IEEE}%
\thanks{Runze Han and Tianyang Zhao are with the School of Electrical
Engineering, Xi'an Jiaotong University, Xi'an 710049, China (e-mail:
zhao@xjtu.edu.cn).}%
\thanks{Peng Wang is with the School of Electrical and Electronic Engineering,
Nanyang Technological University, Singapore 639798.}}

\begin{document}
\maketitle

\begin{abstract}
Conventional reliability assessment of hybrid AC/DC distribution systems
prescribes the fault-isolation boundary and voltage-source-converter (VSC)
availability before restoration, and therefore neglects subsequent VSC outages
caused by protection--IBR interactions. This paper proposes a transition-aware
reliability assessment method. A pre-restoration mass-matrix
differential-algebraic equation (DAE) first generates the fault-to-isolation
trajectory. For each initiating contingency, a trajectory-class mapping then
partitions operating conditions into mutually exclusive classes according to
the DAE-derived unavailable-VSC set supplied to restoration. The class
frequencies inherit the joint-contingency frequencies and
operating-state probabilities. Conditions in one trajectory class, termed a
restoration-entry state, reuse the same DAE-derived entry, while their original
load, generation, and active/reactive load-shedding MILPs remain state specific.
Thus, transient states are compressed without approximating the restoration
problem. $N$--1, $N$--2, and $N$--3 joint contingencies are considered.
On a modified IEEE 33-bus hybrid AC/DC system, transient modeling increases EENS
from 0.090556 to 0.260831~MWh/yr. The 552 contingency--operating-state conditions
form 92 trajectory classes. Leave-one-state-out validation gives 100\%
class-identification accuracy and zero held-out EENS error, while 100000
projected samples yield 5.76--5.99 times speedup. The results verify the
reliability impact of protection--IBR-induced transitions and the computational
benefit of trajectory-class reuse.
\end{abstract}

\begin{IEEEkeywords}
Distribution system reliability, hybrid AC/DC distribution system,
inverter-based resource, protection, service restoration, transient simulation,
voltage-source converter.
\end{IEEEkeywords}

\section{Introduction}

\subsection{Motivation}

\IEEEPARstart{H}{ybrid} AC/DC distribution systems integrate AC and DC feeders,
inverter-based resources (IBRs), energy storage systems, and interlinking
voltage-source converters (VSCs). The VSCs provide controllable power transfer
between the two subsystems and are therefore important resources for post-fault
service restoration. In conventional reliability assessment, the isolated
fault section and the resources available after isolation are normally given
before a network reconfiguration problem is solved. However, the fault current
of an IBR-dominated system is determined by converter control and current
limitation. The resulting voltage and current trajectories affect relay
operation, fault ride-through (FRT), converter blocking, and converter tripping
\cite{baeckeland2022stationary,azizi2024gridcode,fan2023faultrecovery,
cisneros2024protection}. Consequently, an initiating fault may cause additional
VSC outages before network reconfiguration starts, changing the power-transfer
paths and generation resources available to the restoration controller. A
steady-state restoration model cannot determine these outages by itself, while
performing a differential-algebraic equation (DAE) simulation for every
operating state and every $N$--$k$ contingency results in a heavy computational
burden. It calls for a reliability assessment method that determines the
post-fault restoration entry condition from the transient response and reuses
this information without simplifying the subsequent restoration calculation.

\subsection{Literature Review}

The reliability of active distribution systems has been extensively studied by
analytical and simulation-based approaches. In analytical approaches,
component failure rates and repair times are combined with network topology to
calculate annual reliability indices. Optimization models have been introduced
to represent the detailed placement of circuit breakers and switches and to
determine post-fault network reconfiguration \cite{li2020switchplacement,
li2020reconfiguration}. Distributed energy resources (DERs) and microgrids have
also been included to account for islanded operation and emergency power supply
\cite{karngala2021der}. These models can enforce radiality, power-flow limits,
converter capacities, and load shedding during restoration. Nevertheless, the
fault-isolation boundary and the availability of converters and DERs are
usually treated as known inputs. Thus, the restoration result may be optimistic
when protection--converter interaction removes a resource that is assumed to be
available.

To improve the representation of DERs and microgrids, detailed subsystem
simulation has been combined with analytical reliability assessment. In
\cite{omri2024microgrid}, sequential Monte Carlo simulation and an islanded
resource-management problem are used to derive a multi-state generation model
for each microgrid, which is then incorporated into distribution-system
reliability evaluation. Sequential Monte Carlo methods can also represent
chronological component outages and the operation of decentralized power
systems \cite{breve2024flow}. These studies demonstrate that a detailed model
can be calculated in advance and reused in a large number of reliability
states. However, the retained states mainly describe available generation
capacity or component availability over the scheduling and restoration time
scale. Whether fault clearing and FRT leave an interlinking converter available
at the beginning of restoration is still prescribed rather than obtained from
the coupled transient response.

In parallel, detailed studies have investigated the fault behavior of IBRs and
their interactions with protection systems. Current-limiting control, FRT, and
post-fault recovery strongly affect the current magnitude, voltage support, and
recovery of grid-forming and grid-following converters
\cite{li2022adaptivefrt,qoria2023overcurrent,fan2023faultrecovery,
azizi2024gridcode}. Protection studies further show that the limited and
control-dependent fault current of IBRs can alter relay sensitivity and
coordination \cite{cisneros2024protection}. Time-domain DAE simulation is a
well-established means of representing these control and network interactions
\cite{lara2023psd}, and IEEE Std. 1547 specifies voltage and frequency
ride-through requirements for DERs \cite{ieee1547,ieee1547a}. These transient
studies provide the physical models required to determine converter survival,
but they normally stop after the fault response and do not calculate annual
distribution reliability indices through post-fault network reconfiguration.

Only a few studies have linked the fast protection and converter response to
the inputs of a distribution restoration model. A direct combination is
computationally demanding because the DAE and the restoration mixed-integer
linear programming (MILP) problem must be solved repeatedly for many
contingencies and operating states. The difficulty becomes more significant
for higher-order contingencies, for which the number of simultaneous outage
combinations increases rapidly. More importantly, operating states cannot be
grouped only by load level or by the initiating contingency: two states can
enter restoration with different available VSCs after experiencing different
protection and FRT responses. Therefore, a reusable link between transient
simulation and conventional reliability restoration is still required.

\subsection{Contributions}

To determine the actual post-fault inputs of the restoration problem, a
transition-aware reliability assessment method is proposed in this paper. For
each initiating contingency and operating state, a coupled AC/DC DAE model is
first used to determine the VSCs that become unavailable before restoration.
For each initiating contingency, operating states are progressively partitioned
into trajectory classes according to the DAE-derived unavailable-VSC vector;
each class defines one restoration-entry state. The load and
generation levels are not averaged: the original restoration MILP, including
active and reactive load shedding, is solved for every operating state. The
main contributions of this paper can be summarized as follows:

\begin{enumerate}
    \item A transition-aware reliability assessment model is proposed for
    hybrid AC/DC distribution systems. The coupled network, converter, FRT, and
    protection response is represented by a mass-matrix DAE, and its output is
    converted into the unavailable-VSC set supplied to post-fault restoration.
    Explicit $N$--1, $N$--2, and $N$--3 joint contingencies are included without
    multiplying marginal failure rates or treating converter trips as new
    initiating failures.

    \item A trajectory-class theory is developed to compress the DAE-derived
    restoration entries. An ordered subset tree constructs high-order joint
    events without duplication, while the inverse images of the complete
    DAE-to-restoration interface define the unique coarsest lossless partition
    for each event. The class frequencies conserve initiating-event exposure,
    and the class-conditioned EENS is proved equal to direct state-by-state
    evaluation when each state retains its own load, generation, and restoration
    MILP. The theory separates class compression from the combinatorial
    $N$--$k$ event space and gives an EENS remainder bound for certified order
    truncation. A DAE library and leave-one-state-out test enable reuse without
    claiming to compress the restoration optimization; unresolved frequency is
    retained explicitly.
\end{enumerate}

% The remainder of the paper is organized as follows. Section II defines the
% hierarchical controller architecture, reliability assessment scheme, and
% modeling assumptions. Section III formulates the initiating
% failure exposure, progressive consequence partition, and reliability metrics.
% Section IV presents the adaptive equivalent-state calculation. Section
% V reports the numerical studies, and Section VI concludes.

\section{Problem Statement}
\label{sec:problem}

\begin{figure*}[!t]
\centering
\includegraphics[width=\textwidth]{figures/research_object.pdf}
\caption{Hybrid AC/DC distribution system and transition-aware reliability
assessment. (a) Local VSC controls continuously determine converter currents,
whereas protection functions measure their assigned zones and trip only the
associated breakers or converters. The distribution-level controller receives
the isolated topology and available-resource status after the DAE simulation.
(b) For each joint $N$--$k$ contingency and operating state, the pre-restoration
DAE determines the VSCs that are unavailable after fault clearing. The
restoration MILP is then executed once using the initiating outages, the
DAE-derived VSC outages, and the original operating-state load and generation.
(c) Conventional assessment prescribes the restoration inputs. The proposed
assessment obtains them from protection--IBR interaction and reuses only the
restoration-entry state in subsequent reliability samples.}
\label{fig:research_object}
\end{figure*}

\subsection{Hybrid AC/DC Distribution Systems with Hierarchy Controllers}
\label{sec:research_object}
The assessed system is a hybrid AC/DC distribution network
$G=(\mathcal N^{\rm ac}\cup\mathcal N^{\rm dc},\mathcal E^{\rm ac}\cup
\mathcal E^{\rm dc}\cup\mathcal V)$, where $\mathcal N^{\rm ac}$ and
$\mathcal N^{\rm dc}$ are the AC and DC buses, $\mathcal E^{\rm ac}$ and
$\mathcal E^{\rm dc}$ are the corresponding branches, and $\mathcal V$ is the
set of interlinking VSCs. The system also contains loads $\mathcal D$,
converter-interfaced resources $\mathcal R$, controllable switches
$\mathcal E^{\rm sw}$, and protection devices $\mathcal Z$. The set of
physical components for which initiating outages are studied is denoted by
$\mathcal U\subseteq\mathcal E^{\rm ac}\cup\mathcal E^{\rm dc}\cup
\mathcal V$. As shown in Fig.~\ref{fig:research_object}(a), the physical
network is coordinated by local converter control, protection, and a
distribution-level restoration controller.

For each VSC or converter-interfaced resource $v$, the local controller uses
its terminal voltage, current, active power, and DC-link measurements to form
the current reference and operating mode. Current limitation, active/reactive
current priority, and fault ride-through (FRT) are included in this layer.
Therefore, local VSC control acts continuously during the fault and changes the
AC/DC voltage and current response seen by the protection devices. It does not
directly determine the post-fault switching plan.

Each protection function $z\in\mathcal Z$ is assigned a measurement set
$\mathcal M_z$, a protection zone $\mathcal P_z$, and a trip set
$\mathcal B_z$. Its pickup and definite-time states are driven only by the
available relay measurements. Once its operating condition is satisfied, the
corresponding breaker or converter in $\mathcal B_z$ is opened and the AC/DC
network equations are reinitialized. Hence, converter control and protection
are coupled through the same bus voltages and branch currents, although the two
controllers do not exchange their internal states. Complete observation of all
physical and controller states is not assumed.

The distribution-level controller is separated from the preceding transient
response. After fault isolation and completion of the DAE simulation, it
receives the isolated components, breaker positions, and reported VSC and DER
availability, and then solves the network reconfiguration and load-restoration
problem once. The resulting switch operations and power schedules are not fed
back into the DAE trajectory considered in this paper. Thus, local converter
and protection controls determine the actual restoration input, while the
system-level controller determines the load that can be supplied from this
input under the AC/DC network constraints.

\subsection{Reliability Assessment Scheme and Assumptions}
\label{sec:assessment_scheme}
Figure~\ref{fig:research_object}(b) gives the complete calculation procedure.
For an initiating contingency $k$ and operating state $\omega$, a consistent
pre-fault power-flow solution is first obtained. The initiating fault is then
applied to the coupled network, VSC-control, FRT, and protection DAE. When the
fault is cleared and all modeled protection actions are completed, the DAE
provides the set of unavailable VSCs at the restoration instant. This set,
together with the initiating outages, is supplied to the conventional
restoration MILP. The original load and generation of $\omega$ are retained,
and active and reactive load shedding are optimized for that operating state.
The annual reliability indices are finally obtained from the contingency
frequency, operating-state probability, restoration duration, and curtailed
load. The proposed method therefore changes the restoration-entry information,
but does not replace the conventional restoration model.

Four assumptions are made in this paper. First, $N$--1, $N$--2, and $N$--3
contingencies are represented as explicitly authored simultaneous joint events;
their frequencies are input data and are not calculated by multiplying
marginal component failure rates. Second, the pre-restoration model is a
balanced positive-sequence, phasor-domain mass-matrix DAE with average-value
VSCs, explicit current limits, FRT, and definite-time protection; traveling
waves, breaker arcs, and switching harmonics are outside the present scope.
Third, the DAE is terminated at a specified restoration instant after fault
isolation, and the system-level restoration is executed once without feedback
to the preceding trajectory. Fourth, a failed DAE initialization, event-time
network solution, or restoration optimization is retained as unresolved; no
normal restoration state or zero load shedding is assigned to an unresolved
condition.

\input{section_iii}
\input{section_iv}

\input{case_studies}

\section{Conclusion}
In this paper, a transition-aware reliability assessment method for hybrid
AC/DC distribution systems with IBRs is proposed. A pre-restoration DAE model
is used to determine whether the initiating fault causes additional VSC
outages through converter control, FRT, and protection interaction. The
resulting unavailable-VSC vector is supplied to the conventional restoration
MILP. For each initiating event, a terminal-status map progressively partitions
the DAE trajectories into finite, mutually exclusive trajectory classes. This
partition conserves the initiating-event frequency and the state-specific EENS.
Operating states in one class share only the DAE-derived restoration entry;
their original load, generation, and active/reactive load-shedding calculations
are retained. The method therefore improves the physical basis of the
restoration input and reduces repeated DAE calculations without compressing the
restoration problem.

The modified IEEE 33-bus results demonstrate that prescribing VSC availability
can considerably underestimate EENS under the stated $N$--1, $N$--2, and
$N$--3 design frequencies. The leave-one-out tests verify trajectory-class
identification within the six evaluated operating states, and the complete
timing model demonstrates a net computational saving, although the specified
20-times speedup target is not reached. The modified IEEE 123-bus study further
shows that numerical closure of the DAE library is a necessary condition for
reporting feeder-wide indices: 158 unresolved transient conditions prevent its
resolved-submeasure EENS and timing projection from being treated as complete
validation. Future work will improve the event-time AC/DC algebraic solution,
use utility-calibrated joint-contingency and operating-state data, and validate
the restoration entry against detailed three-phase and vendor-specific
converter and protection models.

\bibliographystyle{IEEEtran}
\bibliography{references}

\end{document}
